\section{CSS_Intro}\label{sec:CSS_Intro}

\begin{frame}
    \frametitle{Estilo, CSS}
    
    Si entraron a cualquier sitio web de los últimos 20 años, seguramente hayan notado 
    que se ven bastante más complejos que lo que hicimos hasta ahora con HTML.

    Entra CSS, Cascading Style Sheets. El lenguaje existe para cambiar la 
    forma en la que se ven los documentos HTML. Hasta ahora dependíamos de 
    cómo los distintos elementos fueran renderizados por el navegador, pero 
    ahora podemos cambiar eso completamente.

\end{frame}

\begin{frame}
    \frametitle{Quién necesita JS?}

    Como dato curioso, HTML y CSS juntos son turing complete. Eso significa que, 
    técnicamente, se pueden usar para hacer cualquier cosa que los otros lenguajes 
    de programación puedan hacer, incluyendo JavaScript. Esto no significa que 
    eso sea una buena idea. Estos lenguajes no fueron pensados para eso y tratar 
    de usarlos de esa forma sería increíblemente incómodo y no valdría la pena.

    De todas formas, cuando alguien les diga que HTML y CSS no son lenguajes de 
    programación, pueden responder con este dato.

\end{frame}

\begin{frame}[fragile]
    \frametitle{Sintaxis de CSS}
    
    CSS funciona con \emph{reglas}. Un archivo .css (o una tag \inlinecode{style} 
    en un archivo .html) contiene tantas reglas como uno quiera, una atrás de 
    otra, definiendo a qué partes del documento se le deberían aplicar qué 
    propiedades. Estas reglas siguen el siguiente formato:

    \begin{lstlisting}
selector {
    propiedad1: valor1;
    propiedad2: valor2;
    ...
}
    \end{lstlisting}

    El selector indica qué elementos van a ser afectados por la regla. Las 
    propiedades indican cómo van a ser afectados.

\end{frame}

\begin{frame}[fragile]
    \frametitle{Selectores, clases e ids}
    
    Los selectores más básicos son selectores de tipo, de clase y de id. 

    Un selector de tipo dice que la regla va a aplicar a todos los elementos 
    del tipo indicado. Por ejemplo, si el selector es \inlinecode{p}, todos los 
    elementos \inlinecode{p} van a tener la regla.

    Los selectores de clase y de id funcionan igual, pero hay que marcarlos 
    con un punto y con un numeral respectivamente. Por ejemplo, si el 
    selector es \inlinecode{.contenedor}, todos los elementos con la clase 
    ``contenedor'' van a tener esa regla. Lo mismo con \inlinecode{#titulo-principal} 
    y el elemento con id ``titulo-principal''.

\end{frame}

\begin{frame}[fragile]
    \frametitle{Clases e ids?}
    
    Pero, qué son las clases y los ids?

    Todo elemento de HTML puede tener los atributos \inlinecode{class} e 
    \inlinecode{id}. El primero es una lista de nombres separados con espacios. 
    El segundo es un solo nombre.

    La idea de las clases es que sean indicadores de características compartidas 
    por múltiples elementos a lo largo del documento. También, un elemento puede 
    (y suele) tener múltiples clases. La diferencia con los ids es que el id debe 
    ser único. 

    El motivo por el que estos atributos existen es para poder identificar ciertos 
    elementos o conjuntos de elementos del sitio. Como ya vimos, en CSS esto nos 
    deja aplicar reglas de forma mucho más específica. En el caso de JS, nos 
    permitiría encontrar e interactuar elementos de la misma forma.

\end{frame}

\begin{frame}[fragile]
    \frametitle{Ejemplo de HTML y CSS usando clases e ids}
    
\begin{lstlisting}
<div class="contenedor" id="contenedor-principal">
    <h1 id="titulo-principal">Buenas tardes<h1>
</div>
<div class="contenedor azul">...</div>
\end{lstlisting}

\begin{lstlisting}
.contenedor {
    display: flex;
}
#contenedor-principal {
    background-color: orange;
}
#titulo-principal {
    color: red;
}
.azul {
    background-color: blue;
}
\end{lstlisting}

\end{frame}

\begin{frame}[fragile]
    \frametitle{Más selectores}
    
    Pero esos tres no son los únicos tipos de selectores. Se pueden combinar 
    entre sí de distintas formas para armar selectores mucho más complejos 
    que nos permiten ser mucho más específicos sobre a qué partes del sitio 
    se les debería aplicar qué reglas.

    Algunos ejemplos son:
    
    \begin{itemize}
        \item \inlinecode{A, B} para indicar que la regla debe aplicar a 
              \inlinecode{A} y a \inlinecode{B}.
        \item \inlinecode{A B} para indicar que la regla debe aplicar a los 
              \inlinecode{B} que sean descendientes de \inlinecode{A}.
        \item \inlinecode{A:PC} para indicar que la regla debe 
              aplicar a \inlinecode{A} cuando tiene la pseudoclase \inlinecode{PS}.
    \end{itemize}

    Entre otros.

\end{frame}

\begin{frame}[fragile]
    \frametitle{Valores para las propiedades}
    
    Ya hablamos de los selectores. De las propiedades no hay mucho que decir, 
    hay un montón, cada una tiene un nombre único y ciertos valores que puede 
    tomar. Pero esos valores no son ni un poquito sencillos.

    Por ahora vamos a concentrarnos en dos tipos de valores posibles:

    \begin{itemize}
        \item \emph{keyword}: Son valores con nombres específicos y hacen 
              cosas específicas dependiendo de para qué propiedad están siendo usados.
        \item Magnitudes: Son valores de la forma $Nu$, donde $u$ es la unidad y 
              $N$ es la cantidad de dicha unidad. Estas unidades pueden ser 
              absolutas o relativas. Las absolutas son siempre lo mismo, sin importar 
              en donde aparezcan$^*$. Las relativas representan distintas magnitudes 
              dependiendo del contexto, ya sea dado por el elemento al que se está 
              aplicando la regla o por el propio navegador en el que está viéndose 
              el sitio.
    \end{itemize}

    Más adelante vamos a profundizar un poco sobre las unidades, pero por ahora 
    algunas útiles e interesantes son \inlinecode{px}, \inlinecode{cm}, \inlinecode{\%} 
    y \inlinecode{em}.

\end{frame}

\begin{frame}[fragile]
    \frametitle{Algunas propiedades básicas}
    
    Algunas propiedades básicas para que empecemos a tocar un poco el estilo 
    de nuestros sitios son:

    \begin{itemize}
        \item \inlinecode{color} para cambiar el color del texto. El color puede 
              escribirse por nombre, por código hexadecimal, o con una variable.
        \item \inlinecode{background-color} para cambiar el color del fondo. Recibe 
              lo mismo que \inlinecode{color}.
        \item \inlinecode{border} para ponerle un borde al elemento. Recibe un grosor, 
              un estilo y un color.
        \item \inlinecode{font-size} para cambiar el tamaño del texto. Recibe una 
              magnitud.
        \item \inlinecode{font-family} para cambiar la fuente del texto. Recibe el 
              nombre de la fuente (o familia de fuentes).
        \item \inlinecode{background-image} para poner una imagen de fondo. Recibe 
              una dirección a la imagen o algunos otros valores medio raros.
        \item \inlinecode{text-decoration} para agregarle algunos detalles al texto, 
              como por ejemplo subrayado. Recibe valores concretos como \inlinecode{underline}.
    \end{itemize}

\end{frame}

\begin{frame}[fragile]
    \frametitle{Ejercicio - Historia}
    
    En este ejercicio tienen que retocar y arreglar el estilo del documento 
    sobre la historia de nuestra carrera. Pueden usar las propiedades que vimos 
    en la diapositiva anterior y también pueden buscar otras propiedades que 
    les interesen en internet.

    Este es un buen momento para mencionar 
    \href{https://developer.mozilla.org/}{developer.mozilla.org}, lugar en el que 
    está documentado casi todo lo relacionado a HTML, CSS y JS. El sitio tiene 
    buenas explicaciones, tutoriales, ejemplos y ejercicios. Está buenísimo.
    Pueden buscar ahí lo que quieran para resolver los ejercicios del taller.

    Eso si, por ahora eviten meterse con reglas que cambien el ``layout'' del sitio. 
    Eso incluye padding, márgenes, posicionamiento, display, etc. Después de este 
    ejercicio nos vamos a meter con esos temas.

\end{frame}